%==============================================================================
%  05-organizational-readiness.tex — Section 4: Organizational Readiness
%==============================================================================
\strategyPart{11 \textbar\ Organizational Readiness}{Maturity Levels and Organizational ``Stocks''}

AI readiness is frequently mistaken for technical readiness. It is not. MIT
CISR's research on AI-enabled organizations shows that the binding constraints
on AI value creation are organizational: strategy, data, talent, governance
and culture [MIT CISR]. These constraints behave like stocks: they accumulate
slowly, compound over time, and cannot be rented at the point of need.

\section{The three organizational stocks}

\vspace{-0.3em}
\aiCard{%
  \aiCardTitle{Human Capital}\\[0.4em]
  AI-literate leadership, data and ML talent, and a workforce trained to work
  alongside machines.
}\par\vspace{0.7em}

\aiCard{%
  \aiCardTitle{Data Capital}\\[0.4em]
  Trusted, governed, accessible data. Poor data quality costs the average
  organization \$12.9~million per year [Gartner 2021].
}\par\vspace{0.7em}

\aiCard{%
  \aiCardTitle{Organizational Capital}\\[0.4em]
  Redesigned workflows, aligned incentives, decision rights and cross-functional
  collaboration.
}\par

\section{The MIT CISR readiness lens}

MIT CISR research on the AI-enabled organization emphasizes that readiness is
the alignment of these stocks with strategy: organizations that treat AI as an
isolated IT program accumulate tools without accumulating capability
[MIT CISR]. The empirical evidence supports this framing. In the MIT Sloan--BCG
global study, 90\% of companies have invested in AI, but fewer than two in five
report business gains, and 40\% of significant investors report no bottom-line
impact [MIT SMR--BCG 2019]. The minority who succeed — the ``pioneers'' —
combine strategy, organizational behavior and technology in equal measure.

\section{Cultural readiness: humans and machines}

HBR's landmark study of 1,500 companies concluded that the greatest performance
gains occur when humans and machines work together, not when machines replace
humans [HBR, Wilson \& Daugherty 2018]. The ``collaborative intelligence''
model requires cultural readiness on three fronts:
\begin{itemize}
  \item \textbf{Trust} — employees must trust the data and the model's judgment.
  \item \textbf{Augmentation, not substitution} — roles are redesigned around
        human--machine teams, protecting employment security to reduce
        resistance.
  \item \textbf{Literacy} — every manager, not only data scientists, must read
        and interrogate AI outputs.
\end{itemize}

\section{Maturity stages for each stock}

\begin{center}
  \renewcommand{\arraystretch}{1.45}
  \rowcolors{2}{inSnow}{white}
  \fitwidth{\begin{tabular}{>{\raggedright\arraybackslash}p{3.2cm}>{\raggedright\arraybackslash}p{3.2cm}>{\raggedright\arraybackslash}p{5.2cm}}
    \rowcolor{inNavy}
    \hcell{Stage} & \hcell{Human / Data / Org capital} & \hcell{Evidence of readiness} \\
    \textbf{Foundations} & Sparse & Isolated proofs, data silos, no AI governance. \\
    \textbf{Capable} & Developing & Repeatable delivery, governed data, trained teams. \\
    \textbf{Differentiating} & Strong & AI embedded in core workflows with measured P\&L impact. \\
    \textbf{Reinvented} & System-wide & Business model, products and culture are AI-native. \\
  \end{tabular}}
\end{center}

The executive question is not ``how many models do we run'' but ``how mature
is each stock'' — because the stock that is weakest determines the ceiling on
value.
